\documentclass[a4paper,11pt]{ltjsarticle}
\usepackage[haranoaji]{luatexja-preset}
\usepackage{amsmath,amssymb}
\usepackage[top=12mm,bottom=10mm,left=14mm,right=14mm]{geometry}
\usepackage{array}
\usepackage{tikz}
\usetikzlibrary{calc}
\pagestyle{empty}
\setlength{\parindent}{0pt}
\newlength{\qcol}\setlength{\qcol}{0.47\textwidth}
\newlength{\qgap}
\newlength{\anscolw}\setlength{\anscolw}{0.30\textwidth}
\newlength{\ansh}
\newcommand{\Q}[2]{\parbox[t]{\qcol}{\raggedright (#1)\hspace{0.6em}$\displaystyle #2$}}
\newcommand{\QR}[4]{\Q{#1}{#2}\hfill\Q{#3}{#4}\par\vspace{\qgap}}
\newcommand{\anscell}[1]{\rule[-\ansdrop]{0pt}{\ansh}(#1)}
\newlength{\ansdrop}

\begin{document}
{\large\bfseries 計算問題　26}\hfill 名前(\underline{\hspace{32mm}})\par\vspace{3mm}
■（1）〜（16）の計算をしなさい。（17）〜（21）は方程式を解きなさい。\par\vspace{3mm}
\setlength{\qgap}{5.4mm}
\QR{1}{-2^2+(-4)^2-4}{2}{(-6)\times(-6)}
\QR{3}{(8a-16)\div4}{4}{2(4x-2)-(4x-9)}
\QR{5}{-0.2-(-4.4)-1.9}{6}{\dfrac{x-2}{4}\times16}
\QR{7}{(4x-2)-(9x+4)}{8}{(-9)+(-7)-(-9)}
\QR{9}{(-8)\div(+4)}{10}{10-(-4)\div\left(-\dfrac{4}{4}\right)}
\QR{11}{\dfrac{2}{4}(4x-36)-\dfrac{4}{9}(9x+18)}{12}{3\times\{4-(-4+1)\times2\}}
\QR{13}{2\div4\times(-16)}{14}{(2a+4)+(4a-9)}
\QR{15}{\dfrac{2}{4}\div4+\dfrac{9}{4}}{16}{\dfrac{1}{2}(8x-8)+\dfrac{1}{9}(36x+36)}
\QR{17}{2x-4=-4}{18}{2(x+4)=3x+6}
\QR{19}{0.2x+0.4=0.6}{20}{\dfrac{3}{5}x+2=\dfrac{1}{10}x-1}
\vspace{1mm}
\Q{21}{\dfrac{2x-4}{4}-\dfrac{x+9}{4}=1}\par\vspace{3mm}
\setlength{\ansh}{9.5mm}\setlength{\ansdrop}{6mm}%
\begin{center}\begin{tabular}{|p{\anscolw}|p{\anscolw}|p{\anscolw}|}\hline
\anscell{1} & \anscell{2} & \anscell{3} \\\hline
\anscell{4} & \anscell{5} & \anscell{6} \\\hline
\anscell{7} & \anscell{8} & \anscell{9} \\\hline
\anscell{10} & \anscell{11} & \anscell{12} \\\hline
\anscell{13} & \anscell{14} & \anscell{15} \\\hline
\anscell{16} & \anscell{17} & \anscell{18} \\\hline
\anscell{19} & \anscell{20} & \anscell{21} \\\hline
\end{tabular}\end{center}

\end{document}
